\bookchapter{第六十七章：打磨同發佈——測試、調試同上架 App Store}

\chapterintro{App 整好咗，跟住點呢？呢一章回答“由行得到”到“交到貨”嘅三件事：點樣高效咁調試搵 bug、點樣寫測試防回歸、點樣將 App 送上 App Store 畀全世界下載。上架唔止係儀式感——佢係將一件作品變成產品嘅分水嶺。}

\booksection{一、調試三板斧：print、斷點、預覽}

\begin{center}
\begin{tabularx}{\linewidth}{|X|X|X|}
\hline
\textbf{工具} & \textbf{做乜嘢} & \textbf{點樣用} \tabularnewline
\hline
print + 控制檯 & 睇運行日誌 & 代碼入面 print("值=\textbackslash{}(x)")，喺 Xcode 底部控制台睇輸出 \tabularnewline
斷點調試 & 令程式停喺某一行，睇變量值 & 喺行號旁邊㩒一下出藍標，運行到嗰度就停，逐步行 \tabularnewline
SwiftUI 預覽 + 調試器 & 即時睇界面同視圖層級 & \#Preview 畫布 + 運行時嘅 View Hierarchy 檢查器 \tabularnewline
\hline
\end{tabularx}
\end{center}

\begin{center}\small\color{gray}調試三板斧\end{center}

\begin{tcbinfo}{小貼士：崩潰咗點樣讀？}
Swift 崩潰嗰陣控制台會打印 “Fatal error” 加原因，Xcode 會停喺出錯嗰一行並高亮。真機上面嘅歷史崩潰可以喺 Xcode Window Organizer Crashes 入面睇（需要付費賬號）。關鍵詞優先搜“錯誤原因第一句 + swiftui”——精確嘅問題，先至有精確嘅答案。

\end{tcbinfo}

\booksection{二、測試：畀代碼上保險}

手測靠唔住——改咗呢個功能，另一個靜靜哋壞咗。測試就係“後悔藥”兼“防彈衣”。iOS 用 XCTest 框架，分兩種：

\begin{center}
\begin{tabularx}{\linewidth}{|X|X|X|X|}
\hline
\textbf{層級} & \textbf{測乜嘢} & \textbf{Target} & \textbf{速度} \tabularnewline
\hline
單元測試 & 純邏輯（計算、校驗） & MyFirstAppTests & 秒級 \tabularnewline
UI 測試 & 真實界面行為 & MyFirstAppUITests & 分鐘級 \tabularnewline
手動測試 & 真人感受、體驗 & —— & 永遠需要 \tabularnewline
\hline
\end{tabularx}
\end{center}

\begin{center}\small\color{gray}iOS 測試層級（模板項目已預置兩個 Test Target）\end{center}

\begin{lstlisting}[language=Swift]
import XCTest

// 把"金額字串轉數字"呢種邏輯抽成純函數
func parseAmount(_ input: String) -> Double? {
    guard let v = Double(input.trimmingCharacters(in: .whitespaces)),
          v > 0 else { return nil }
    return v
}

final class AmountTests: XCTestCase {
    func testEmptyInputReturnsNil() {
        XCTAssertNil(parseAmount(""))
    }
    func testNegativeReturnsNil() {
        XCTAssertNil(parseAmount("-5"))
    }
    func testNormalValueParses() {
        XCTAssertEqual(parseAmount("12.5"), 12.5)
    }
}
\end{lstlisting}

行測試：CmdU。全綠之後再改代碼，心裏有底。記賬本入面嘅金額校驗、匯總計算，都應該有單元測試罩住。

\booksection{三、簽名同打包：畀 App 蓋上“印章”}

每個 iOS App 都必須簽名先安裝到——證明“呢個包出自你手，冇畀人改過”。層級關係：

\begin{itemize}
  \item 學習階段：免費 Apple ID + Xcode 自動簽名（Personal Team），真機行得到，7 日續期一次；
  \item 分發階段：付費開發者賬號（99 美元／年）簽發證書同描述檔，先可以上架同長期安裝。
\end{itemize}

打包：Xcode 菜單 Product Archive（歸檔），構建完成之後喺 Organizer 入面 Distribute App 揀 App Store Connect 上傳。歸檔必須用 Generic iOS Device 做目標（唔係模擬器），呢一步係新手成日唔記得嘅。

\textbf{避坑提醒：證書丟咗 = 更新斷纜}\par

簽名證書同私鑰一旦丟失，你就冇得再畀舊 App 發更新——用戶只可以裝“全新”應用。妥當備份證書（鑰匙串導出 .p12）並記好密碼。好在蘋果提供咗“自動簽名”同雲端證書管理（Xcode 入面剔選 Automatically manage signing），可以慳返大部分手工操作嘅坑。

\booksection{四、上架 App Store 全流程}

註冊開發者賬號

準備商店素材

上傳構建包

填隱私與審覈資訊

蘋果審覈（1-7 天）

上架

\begin{center}\small\color{gray}圖 5：App Store 上架六步走\end{center}

\begin{itemize}
  \item 註冊：developer.apple.com，個人賬號 99 美元／年；
  \item 素材：圖標（1024×1024，咪帶圓角——蘋果自己切）、截圖（iPhone 各尺寸）、描述、私隱政策連結；
  \item 私隱問卷：如實申報收集咗哪些數據——唔上網嘅記賬本可以"唔收集任何數據"；
  \item 審核：首次提交通常要等 1-7 日，被拒就按"審核備註"逐條改——蘋果嘅拒絕理由通常寫得好具體，照做就得；
  \item TestFlight：上架之前先將測試版傳畀朋友內測（最多 10000 人），收集一輪反饋再正式提交，專業範十足。
\end{itemize}

\booksection{五、商店素材點樣做：第一印象就係下載量}

商店頁係用戶“行街”嗰陣睇到嘅櫥窗。三個素材決定第一印象：

\begin{itemize}
  \item 圖標：1024×1024，咪用截圖求其頂住。冇設計能力就用 AI 工具生成一個簡潔圖形，喺桌面縮細之後仲一眼認得出先算合格；
  \item 截圖：至少 2 套（iPhone 主流尺寸），展示最核心嘅功能場景；每張配一句貼地文案（"一秒記一筆"好過"高效的記賬解決方案"十倍）；
  \item 描述：第一句講清楚"呢個係乜嘢 + 畀邊個用"，咪堆砌"最專業、最強"——真誠好過誇張。
\end{itemize}

\begin{tcbpractice}{嚟動手：將記賬本"交付"一次}
\booksection{六、崩潰上報同性能優化小清單}

App 上架之前，呢兩件事值得各花半個鐘：

\booksubsection{崩潰監控}

你測試嗰陣冇崩，唔代表用戶部手機唔崩。蘋果喺 Xcode Organizer 入面提供免費嘅崩潰報告（付費賬號可以睇真機崩潰日誌）；更省心嘅方案係接入第三方 SDK（好似 Firebase Crashlytics 咁）。冇崩潰監控就上架，等於合埋眼揸車。

\booksubsection{性能自查四條}

\begin{enumerate}
  \item 列表用 LazyVStack / List：長列表咪用普通 VStack 硬排——懶加載係性能第一課；
  \item 圖片咪硬塞原圖：用 AsyncImage（遠程圖）同系統快取，5MB 原圖足以卡爆舊機型；
  \item 咪喺主線程做重活：網絡、大文件都行 async/await；
  \item 用 Instruments 度一度：Xcode 自帶嘅性能分析器（Time Profiler、記憶體泄漏檢查）——先測量，再優化，咪靠感覺。
\end{enumerate}

(1) 畀 parseAmount 等邏輯補 3 個單元測試（CmdU 行到全綠）；(2) 用斷點親手捉一個自己整出嚟嘅崩潰並整好佢；(3) Product Archive 打一個包，行 TestFlight 邀請朋友內測；(4) 有餘力就註冊開發者賬號，將記賬本真正提交上架。頭三步今晚就做得完。

\end{tcbpractice}

\begin{tcbsummary}{本章小結}
\begin{itemize}
  \item 調試三板斧：print 日誌、斷點、SwiftUI 預覽同視圖檢查器
  \item XCTest：單元測試（CmdU）+ UI 測試 + 手動測試，模板已預置
  \item 簽名：免費賬號自動簽名（7 日續期）；付費賬號先可以分發
  \item 打包：Generic iOS Device + Archive + Distribute；證書務必備份
  \item 上架六步：註冊 素材 上傳 私隱問卷 審核 上架；TestFlight 先內測
\end{itemize}

\end{tcbsummary}
